\documentclass[11pt, a4paper]{article}

% ── Encoding & language ───────────────────────────────────────────────────────
\usepackage[T1]{fontenc}
\usepackage[utf8]{inputenc}
\usepackage[british]{babel}

% ── Page geometry ─────────────────────────────────────────────────────────────
\usepackage[
  a4paper,
  top=25mm, bottom=25mm,
  left=25mm, right=25mm,
  headheight=14pt
]{geometry}

% ── Fonts ─────────────────────────────────────────────────────────────────────
\usepackage{mathptmx}
\usepackage{microtype}

% ── Colour ────────────────────────────────────────────────────────────────────
\usepackage[dvipsnames, table]{xcolor}
\definecolor{headblue}{HTML}{1F4E79}
\definecolor{ruleblue}{HTML}{2E75B6}
\definecolor{accentbg}{HTML}{D6E4F0}
\definecolor{codebg}{HTML}{F5F8FC}
\definecolor{codefg}{HTML}{1F497D}
\definecolor{gray1}{HTML}{595959}
\definecolor{lightgray}{HTML}{CCCCCC}

% ── Section headings ──────────────────────────────────────────────────────────
\usepackage{titlesec}
\titleformat{\section}
  {\color{headblue}\large\bfseries}
  {\color{headblue}\thesection}{1em}{}
  [{\color{ruleblue}\titlerule[1pt]}]
\titleformat{\subsection}
  {\color{ruleblue}\normalsize\bfseries}
  {\color{ruleblue}\thesubsection}{1em}{}
\titleformat{\subsubsection}
  {\color{gray1}\normalsize\bfseries}
  {\color{gray1}\thesubsubsection}{1em}{}
\titlespacing*{\section}{0pt}{18pt}{8pt}
\titlespacing*{\subsection}{0pt}{12pt}{4pt}
\titlespacing*{\subsubsection}{0pt}{8pt}{2pt}

% ── Header & footer ───────────────────────────────────────────────────────────
\usepackage{fancyhdr}
\pagestyle{fancy}
\fancyhf{}
\renewcommand{\headrulewidth}{0.4pt}
\renewcommand{\headrule}{\hbox to\headwidth{%
  \color{ruleblue}\leaders\hrule height \headrulewidth\hfill}}
\fancyhead[L]{\small\color{gray1}FreeFBM Verbalisation}
\fancyhead[R]{\small\color{gray1}Pipeline and Rule Language}
\fancyfoot[C]{\small\color{gray1}\thepage}

% ── Hyperlinks ────────────────────────────────────────────────────────────────
\usepackage[
  colorlinks=true,
  linkcolor=ruleblue,
  urlcolor=ruleblue,
  pdftitle={FreeFBM Verbalisation --- Pipeline and Rule Language v10},
  bookmarks=true,
  bookmarksnumbered=true
]{hyperref}

% ── Tables ────────────────────────────────────────────────────────────────────
\usepackage{booktabs}
\usepackage{array}
\usepackage{tabularx}

% ── Code listings ─────────────────────────────────────────────────────────────
\usepackage{listings}
\lstset{
  basicstyle=\small\ttfamily\color{codefg},
  backgroundcolor=\color{codebg},
  frame=single,
  framerule=0.4pt,
  rulecolor=\color{lightgray},
  breaklines=true,
  keepspaces=true,
  columns=flexible,
  xleftmargin=8pt,
  xrightmargin=8pt,
  aboveskip=8pt,
  belowskip=8pt,
  showstringspaces=false,
  tabsize=2,
}

% ── Note boxes ────────────────────────────────────────────────────────────────
\usepackage[framemethod=default]{mdframed}
\newmdenv[
  backgroundcolor=accentbg,
  linecolor=ruleblue,
  linewidth=1pt,
  innerleftmargin=10pt,
  innerrightmargin=10pt,
  innertopmargin=6pt,
  innerbottommargin=6pt,
  skipabove=8pt,
  skipbelow=8pt,
]{notebox}

% ── Misc ──────────────────────────────────────────────────────────────────────
\usepackage{parskip}
\setlength{\parindent}{0pt}
\setlength{\parskip}{6pt}
\usepackage{enumitem}
\setlist[itemize]{noitemsep, topsep=4pt, leftmargin=1.5em}
\usepackage{amsmath}
\usepackage{amssymb}
\usepackage{needspace}

% ── Convenience macros ────────────────────────────────────────────────────────
\newcommand{\code}[1]{\texttt{\color{codefg}#1}}

% Two-column predicate/helper table
\newenvironment{predtable}{%
  \begin{center}
  \begin{tabular}{>{\leavevmode\ttfamily\color{codefg}\raggedright\arraybackslash}p{0.38\textwidth} p{0.52\textwidth}}
  \toprule
  \normalfont\textbf{Syntax} & \textbf{Meaning} \\
  \midrule
}{%
  \bottomrule
  \end{tabular}
  \end{center}
}
\newcommand{\predrow}[2]{#1 & #2 \\[3pt]}

% ─────────────────────────────────────────────────────────────────────────────
\begin{document}

% ── Cover page ────────────────────────────────────────────────────────────────
\begin{titlepage}
  \vspace*{4cm}
  \begin{center}
    {\color{headblue}\fontsize{32}{38}\selectfont\bfseries
      FreeFBM Verbalisation}\\[0.6em]
    {\color{gray1}\Large
      The Pipeline from Formulas to Sentences}\\[0.2em]
    {\color{gray1}\Large
      and the Rule Language that Drives It}\\[2.5em]
    \textcolor{ruleblue}{\rule{0.6\textwidth}{1.5pt}}\\[0.6em]
    {\color{ruleblue}\large\itshape
      Design Specification\quad\textperiodcentered\quad Version~10}\\[0.6em]
    \textcolor{ruleblue}{\rule{0.6\textwidth}{1.5pt}}\\[3cm]
    {\color{gray1}November 2026\\[0.3em] Describes FreeFBM v1.3.3}
  \end{center}
\end{titlepage}

\tableofcontents
\clearpage

% ═════════════════════════════════════════════════════════════════════════════
\part{The Pipeline}

% ─────────────────────────────────────────────────────────────────────────────
\section{Introduction}

This document describes how FreeFBM turns the formulas of its logic into
sentences in seven natural languages, and the rule language (DSL) in which
the language-specific part of that process is written. Version~10 replaces
the DSL specification v9 (October 2026), which was written after the first
stage of the implementation, when only English existed and the rules were
still JavaScript objects.

v9 already followed the pipeline stage by stage, with the rule language as
one stage among others. That turned out to be the right shape: most design
decisions cut across stages. Whether a variable is introduced as
\textit{some Course} or referred to as \textit{that Course} is decided by a
rule, but can only be settled once the sentence has been linearised; a
negative sentence is a formula rewrite before the rules run, but needs a rule
of its own to be said; a reading's grammatical cases come from the model's
lexicon but are used by the language library when it builds a noun phrase.
This version therefore describes the whole pipeline (Part~I), then the rule
language as a reference that can be read on its own (Part~II), then the
languages and the tools (Parts~III and~IV).

Everything described here runs in FreeFBM v1.3.3 (\code{src/verbalisation/}).
Things that do not exist yet are marked \emph{open} and collected in
Section~\ref{sec:open}.

\subsection{Starting point}

FreeFBM translates every constraint into a formula of a sorted logic with
counting quantifiers, transitive closure and aggregation, FBM logic
(\code{docs/logic/fbm-logic.tex}). The Constraints tab shows these formulas
and their sentences, side by side if wanted. Verbalisation takes a formula as
its \emph{only} input. A named constraint and a hand-written formula
constraint with the same content therefore get the same sentence, and the
correctness of a sentence depends only on the rules that read the formula,
not on a translation of constraint types.

\subsection{Design goals}

\begin{itemize}
  \item \textbf{Formula in, sentence out.} Rules match formula trees; they
        never see constraint types.
  \item \textbf{Total.} Every formula gets a sentence. Each language ends
        every phrase kind with generic rules that accept any node.
  \item \textbf{Idiomatic where possible.} Idiom rules come first and give
        NORMA's wording for the common constraints in English, and wording
        modelled on it in the other languages.
  \item \textbf{Declarative.} Rules are data: patterns, structural guards and
        templates, in rule files (\code{.vr}).
  \item \textbf{Multilingual.} A language is a rule file plus a library; the
        engine, the rewrites, naming and reading selection are shared, and
        one model file carries names and readings in several languages.
  \item \textbf{Traceable and testable.} The engine reports which rules built
        a sentence, and golden files fix the sentence of every formula of a
        reference model, in every language and mode.
\end{itemize}

\subsection{What the pipeline produces}

Two kinds of text are produced through the same machinery:

\begin{itemize}
  \item \textbf{Constraint sentences} (Constraints tab, workbench): one
        sentence per formula, in the positive, negative, default or deontic
        mode (Section~\ref{sec:modes}).
  \item \textbf{Population sentences} (Population tab, Sentences view): one
        sentence per fact or instance, built from readings and preferred
        identifiers with the same reading and noun-phrase helpers
        (Section~\ref{sec:population}).
\end{itemize}

% ─────────────────────────────────────────────────────────────────────────────
\section{Overview}
\label{sec:architecture}

\begin{lstlisting}
FreeFBM model + scope (selection, diagram or whole schema)
    |
    v  [0] formalise()            entries {kind, subject, formula}
    |                             src/logic/constraintFormulas.js
    v  [1] Mode rewrite           negative: phi -> not X       negative.js
    v  [2] Variable naming        the names the formula shows  renderFormula
    v  [3] Normalisation          wildcard rewrite             normalise.js
    v  [4] Rules, per phrase kind sentence / clause / vp       <lang>.vr
    |        - reading selection, lexicon, hyphen binding      engine.js, lexicon.js
    |        - language library: noun phrases, word order,     languages/<lang>.js
    |          agreement, cases
Phrase tree (tokens, groups, choices)
    |
    v  [5] Linearisation          choices, contraction,        engine.js
    |                             spacing, capital, full stop
Sentence (styled tokens + plain text) and a rule trace
\end{lstlisting}

The entry point is \code{verbalise(formula, ctx, language, options)}
(\code{engine.js}). \code{ctx} is the context \code{makeContext(state)}
builds for the formulas (object types, fact types and link fact types by id,
the name of a type as sentences write it, the live language), and
\code{language} a compiled language (Section~\ref{sec:languages}). The
options are \code{negative} (the negative mode) and \code{stats} (a map that
collects how often each rule matched, was rejected by its guard or failed in
its template; Section~\ref{sec:tools}). The result holds the tokens (each
with a style class and a flag for a preceding space), the plain text, the
trace (the names of the rules that built the sentence, outermost first) and
the formula actually said (the rewritten one in the negative mode). The
Constraints tab shows the trace when the pointer rests on a sentence.

Stages [0]--[3] and [5] are shared by all languages; stage [4] is the
language: its rule file and its library. The boundary between the two halves
of stage [4] is a design choice per construction: the rules decide
\emph{what} is said in which order, and the library decides \emph{how} a
phrase is built in that language (Section~\ref{sec:library-interface}).

% ─────────────────────────────────────────────────────────────────────────────
\section{Input: Formulas and Entries}
\label{sec:logic}

\subsection{Entries}

The function \code{formalise} (with the option \code{wholeSchema}) returns the entries of the scope:
each a \code{\{kind, subject, formula\}}, or a note instead of a formula when
a constraint cannot be formalised yet. Besides the explicit constraints there
are entries for what the model says implicitly: subtype disjointness and
participation (implicit participation of non-independent entity types),
implied link fact types of objectified fact types, the uniqueness and
mandatory roles a preferred identifier implies, derivation rules of derived
and semi-derived types and fact types, and formula constraints. The default
entries of Section~\ref{sec:default} are listed in place of the constraint
that is missing. A deontic constraint's formula is wrapped as $O(\varphi)$
(\code{modality.js}).

\subsection{Formula nodes}

The formula is a tree built by \code{src/logic/formula.js}. Variables are
objects compared by identity, each with the type it ranges over (an object
type or an objectified fact type). A variable over the intersection of
several types (the direct supertypes of a derived subtype) lists them all.

\begin{center}
\small
\begin{tabularx}{\textwidth}{>{\ttfamily\color{codefg}}l X}
\toprule
\normalfont\textbf{Node} & \textbf{Meaning} \\
\midrule
forall(vars, body) & $\forall x_1 \ldots x_n\,\varphi$ \\
exists(vars, body) & $\exists x_1 \ldots x_n\,\varphi$ \\
count(op, n, [m,] vars, body) & counting quantifier $\exists^{\le n}$,
  $\exists^{\ge n}$, $\exists^{=n}$, $\exists^{n..m}$ over tuples of
  variables; $\exists^{>n}$ only after the negative rewrite \\
and(args), or(args) & $n$-ary, flattened \\
implies(l, r), iff(l, r), not(a) & connectives \\
true & $\top$ \\
fact(factId, args) & fact type atom; arguments in the fact type's
  \code{roles[]} order; the fact type may be a link fact type
  \code{<id>\_il\_<i>} \\
type(typeId, x) & type membership, $\mathit{Student}(x)$ \\
cmp(op, x, y) & object (in)equality $=$, $\neq$ \\
vcmp(op, s, t) & value comparison $\approx, \not\approx, <, \le, >, \ge$ \\
dpred(op, args) & data predicate (\code{startsWith}, \ldots) \\
tc(star, [v,w], body, [s,t]) & $[\mathsf{TC}_{v,w}\,\varphi](s,t)$, with
  \code{star} the reflexive-transitive closure \\
obligatory(arg) & $O(\varphi)$, a deontic constraint \\
possible(arg) & $\Diamond\varphi$, a default entry (Section~\ref{sec:default}) \\
forbidden(arg) & $F(\chi)$, only after the negative rewrite of $O(\varphi)$ \\
not(arg, mode: impossible) & $\neg\chi$ as the negative rewrite makes it \\
wild(type, v) & a wildcard, only after normalisation (Section~\ref{sec:normalisation}) \\
\bottomrule
\end{tabularx}
\end{center}

Data terms are a variable (its value), \code{const(text, quoted)},
\code{fn(op, args)} for the data operators of \code{dataOperators.js},
\code{ite(cond, then, else)}, and \code{agg(fn, distinct, vars, cond, term)}
for $\mathsf{Count}$, $\mathsf{Sum}$, $\mathsf{Avg}$, $\mathsf{Min}$ and
$\mathsf{Max}$.

% ─────────────────────────────────────────────────────────────────────────────
\section{Presentation Modes}
\label{sec:modes}

A mode is a formula, or a formula rewrite, together with a sentence rule
that reads its top node. The rest of the pipeline does not know about modes.

\subsection{Positive}

The formula as it is: \textit{Each Person has at most one PersonName.}

\subsection{Negative}

As in NORMA, a constraint is said as the situation it rules out
(\code{negative.js}). The formula $\varphi$ is rewritten to $\neg\chi$, where
$\chi$ is the counterexample: $\neg\varphi$ with the negation pushed inward.
$\forall$ becomes $\exists$, $\exists^{\le n}$ becomes $\exists^{>n}$
(\textit{more than}), $\exists^{\ge n}$ becomes $\exists^{\le n-1}$ or a
negated $\exists$, an exact count or a range becomes \textit{fewer or more},
$\wedge$ and $\vee$ swap, $A \to B$ becomes $A \wedge \neg B$, and value
comparisons flip. The result is the node \code{not(}$\chi$\code{, mode:
impossible)}, which every language's sentence rule \code{impossible} reads
as \textit{It is impossible that \ldots}:

\begin{quote}
\textit{It is impossible that some Person has more than one PersonName.}
\end{quote}

A formula without a useful counterexample (a derivation rule, a default
entry) gives no rewrite, and the positive sentence is used, as NORMA does. A
deontic $O(\varphi)$ becomes $F(\chi)$, \textit{It is forbidden that
\ldots}. In negative sentences, uniqueness over roles other than the
subject's is said with \textit{the same} (\code{env.same}): \textit{more than
one Person has the same Email}.

\subsection{Default}
\label{sec:default}

NORMA verbalises what the model allows where a constraint is missing:
\textit{It is possible that more than one Person has the same PersonName.}
FreeFBM's \code{formalise()} gives these entries as $\Diamond\chi$
(\code{possible}), for a missing mandatory role (one that is not explicit,
not identified and not only implicitly mandatory through participation) and
for a missing uniqueness constraint (each maximal set of roles short of all
that contains no uniqueness constraint), and lists each in place of the
missing constraint. The sentence rule \code{possible} reads it.

\subsection{Deontic}

$O(\varphi)$ is said as \textit{It is obligatory that} followed by the
sentence $\varphi$ would be (\code{sentence(?x)}, so its own idiom rules
apply): \textit{It is obligatory that there are exactly 3 Buildings.}

% ─────────────────────────────────────────────────────────────────────────────
\section{Variable Naming}

Variables get exactly the names the formula view shows, computed by
\code{renderFormula} on the formula after the mode rewrite. A variable is
named after its type, with a subscript only when two variables of that type
can be in scope at the same time (\textit{Course}$_1$, \textit{Course}$_2$).
A sentence and its formula therefore always agree. The engine takes only the
subscript from these names; the type name itself comes from the language
(the lexicon entry, Section~\ref{sec:lexicon}), so a Dutch sentence says
\textit{Vak}$_1$ for \textit{Course}$_1$. Wildcards have no subscript. A
variable over several types is named after all of them, joined by \code{\&}.

Type names are written as NORMA's option \textit{Object type name display}
asks (\code{nameStyle.js}): as entered (the default), separated into lower
case words (\textit{InstructionType} $\to$ \textit{instruction type}), or
combined with an upper or lower case lead. A reference mode value type is
written as its two names, the entity type's and the reference mode's.

% ─────────────────────────────────────────────────────────────────────────────
\section{Normalisation}
\label{sec:normalisation}

Normalisation is language-independent and copies the formula; variable
objects are shared with the original so that their names still apply.

\paragraph{Wildcard rewrite.}
A variable $y$ bound by $\exists$ is replaced by a wildcard when it occurs
exactly once in the body, as an argument of a fact atom, and the path from
the $\exists$ to that atom passes only through conjunctions (and nested
$\exists$). The variable is removed from the quantifier, and a quantifier
with no variables left disappears. The rewrite runs bottom-up.
\[
\exists y\,(A \wedge R(x,y)) \;\leadsto\; A \wedge R(x,\_)
\qquad\text{when } y \text{ does not occur in } A.
\]
Disjunctions are not crossed: with an empty population of $y$'s type,
$\exists y\,(A \vee R(y))$ is false while $A \vee \exists y\,R(y)$ need not
be. Negation and the other quantifiers are not crossed either. The rewrite
turns most mandatory, participation and subset formulas into atoms with
\textit{some T} arguments: $\forall x\,\exists y\,R(x,y)$ becomes
$\forall x\,R(x,\_)$, said \textit{Each X R some Y}. A wildcard is written by
the language's \code{wildNP}: \textit{some}, \textit{no} under a negated atom
(\code{env.none}), \textit{any} after \textit{No X} or a negated verb.

% ─────────────────────────────────────────────────────────────────────────────
\section{Readings}
\label{sec:readings}

\subsection{Reading structure}

A reading of an $n$-ary fact type is a role order (a permutation of
$0..n-1$) and $n+1$ text fragments; fragment~$k$ precedes the argument of
role $\mathit{roleOrder}[k]$ and the last fragment follows the last argument.
The primary reading is \code{readingParts} with the identity role order;
further readings are \code{alternativeReadings}. A role order may have
several readings, the first preferred (\code{src/store/readings.js}). Link
fact types of objectified fact types have readings too (\textit{involves}).

\subsection{Per-language readings and names}
\label{sec:lexicon}

The model's lexicon (\code{lexicon.js}) gives each element text per
language. An object type may carry \code{lexicon[lang] = \{name, plural,
gender, forms\}}: gender is the language's own word class (\textit{de}/\textit{het}
in Dutch, \textit{m}/\textit{f} in the Romance languages, five genders in
Polish, \textit{m}/\textit{f}/\textit{n} in German), and \code{forms} the
name in the other cases (Polish, German). A fact type may carry
\code{lexicon[lang].readings = [\{roleOrder, parts, subParts?, pluralParts?,
cases?\}]}:

\begin{itemize}
  \item \code{subParts}: the subordinate-clause word order (Dutch
        \textit{\ldots\ dat Student Vak volgt});
  \item \code{pluralParts}: the reading with a plural subject (\textit{are
        prerequisites for});
  \item \code{cases}: the grammatical case of each slot, in reading order
        (Polish \code{['nom', 'loc']} for \textit{Sala znajduje się w
        Budynku}).
\end{itemize}

The live language, the one the element's own name and readings are in (the
view language, \code{ctx.\allowbreak liveLanguage}), uses the element's own text; its
lexicon entries only add the extra forms. For another language the
lexicon's readings replace the element's own, and without any the element's
own are used. The model file is saved in its base language
(\code{src/store/viewLanguage.js}).

\subsection{Selection}

Selection happens per atom occurrence, when a rule renders an atom
(\code{reading()} in \code{engine.js}):

\begin{itemize}
  \item \textbf{With a subject} (verb phrases): only a reading whose role
        order starts with the subject's role, with nothing before it, is
        used. The subject's slot is left out; a second occurrence of the
        subject is the language's reflexive (\textit{itself}). Without such
        a reading the verb phrase fails and the caller falls back to a
        clause.
  \item \textbf{Without a subject} (clauses): the first reading that starts
        with a variable not filled by an override is preferred (\textit{that
        Course has some CourseTitle}), otherwise the primary reading.
\end{itemize}

The language then shapes the chosen reading for the environment
(\code{lang.shapeReading}): the plural form with a plural subject, an
inflected first verb (\code{lang.verbForm}: negated, plural), the
subordinate or inverted word order in Dutch and German. Shaping may fail,
which fails the reading.

\subsection{Overrides}

A rule may fill a variable's slot with a phrase of its own: a quantified
noun phrase (\textit{at most one Email}), a noun phrase with a relative
clause (\textit{some Offering that involves some Course}), or nothing (the
gap of an object relative clause, \textit{that some AcademicYear starts
in}). In languages with cases an override is a function from the slot's
case to the phrase, so the noun phrase is built in the case the reading
gives the slot (\textit{jakiegoś Przedmiotu}, \textit{einen Kurs}).

\subsection{Hyphen binding}

As in NORMA, a reading may bind words to a role player with hyphens
(\code{hyphenBinding.js}), so that the quantifier goes in front of them:
with \textit{\ldots\ has first- given \ldots}, \textit{each Person has at
most one first given Name}. A word ending in a hyphen binds itself and the
words after it up to the player (left binding); a word starting with a hyphen
binds the words from the player up to itself (right binding,
\textit{au plus un Nom principal}); a doubled hyphen joins without a space
(\textit{ex-Partner}). The hyphens are markers only: diagrams show the
reading as typed, sentences and formulas without them. A reading that binds
words to its first player does not count as starting with it.

% ─────────────────────────────────────────────────────────────────────────────
\section{Rules}
\label{sec:rules}

Stage [4] applies the language's rules. Part~II gives the rule language in
full; this section gives the model of computation.

\subsection{Phrase kinds}

Every rule belongs to one phrase kind. To build a phrase of a kind for a
node, the engine tries the rules of that kind in file order and uses the
first whose pattern matches, whose guard holds and whose template succeeds.

\begin{center}
\begin{tabularx}{\textwidth}{>{\bfseries}l X l}
\toprule
Kind & Produces & Example \\
\midrule
sentence & a whole formula; linearisation adds the capital and full stop &
  \textit{Each Person has some PersonName} \\
clause & a formula inside a sentence; introduced variables are referred
  to as \textit{that X} & \textit{some Room is in that Building} \\
vp & a verb phrase about a subject variable, whose slot is left out &
  \textit{has at most one Email} \\
\bottomrule
\end{tabularx}
\end{center}

Noun phrases are not a phrase kind: they are built by library functions
(\code{quantNP}, \code{ref}, \code{np} in the languages with cases), because
their shape depends on the slot they fill.

\subsection{Failure and totality}

A template \emph{fails} when a sub-phrase it needs cannot be built (a verb
phrase for a subject that has no reading starting with it, a reading the
language cannot shape). The engine then tries the next rule. Failure is
ordinary control flow, and a failed attempt leaves nothing in the trace. The
last rules of each kind are generic: the clause kind has a rule per node type
and a final catch-all, and the sentence kind falls back to a clause. Every
formula therefore gets a sentence, though not always an idiomatic one; the
coverage tool lists the sentences that needed the catch-all
(Section~\ref{sec:tools}).

\subsection{Environment}
\label{sec:env}

Each call carries an environment that templates extend with
\code{with(key: value) \{ \ldots \}} and that rules, guards and the library
read. The keys in use:

\begin{center}
\small
\begin{tabularx}{\textwidth}{>{\ttfamily\color{codefg}}l X}
\toprule
\normalfont\textbf{Key} & \textbf{Meaning} \\
\midrule
subject & the subject variable of a verb phrase (\code{?subj} in vp rules) \\
verb & the form of the subject's verb: \code{plural}, \code{negated}
  (Romance also \code{ne}); carried over to verb phrases about the same
  subject only \\
intro & the variables an \textit{If \ldots\ then \ldots} sentence is about:
  each is introduced at its first mention as a wildcard is and referred back
  to after it (Section~\ref{sec:choices}) \\
refStyle & \code{bare}: references are bare names (derivation rules,
  \textit{Student is a TeachingAssistant if and only if Student assists in
  some Course}) \\
same & variables written \textit{the same X} (negative uniqueness) \\
none & a wildcard is \textit{no X} (and in Romance the verb is negated) \\
negative & set under \textit{No X}: wildcards read \textit{any} \\
form & word order in Dutch and German: main (unset), \code{sub}
  (verb last), \code{inv} (verb first) \\
\bottomrule
\end{tabularx}
\end{center}

A clause starts without a subject and verb form, a sentence also without
\code{subject} and \code{verb}; every other key is inherited.

\subsection{References and wildcards}

Templates never write a variable's determiner themselves; they ask for a
reference (\code{ref(?x)}, or a slot of a reading) and the language writes
it: \code{refNP} for an introduced variable (\textit{that Course},
\textit{dat Vak}, \textit{diesen Kurs}, \textit{the same Email}, or bare),
\code{wildNP} for a wildcard (\textit{some Course}, \textit{een Vak},
\textit{jakiegoś Przedmiotu}). This keeps agreement, gender and case in the
library, and lets the environment change all references of a sub-phrase at
once.

\subsection{First and later mentions}
\label{sec:choices}

In an \textit{If \ldots\ then \ldots} sentence the universally quantified
variables are introduced in the condition. NORMA writes each with
\textit{some} at its first mention and \textit{that} at every later one:

\begin{quote}
\textit{If some Staff member teaches some Offering that involves some Course
then that Staff member is qualified to teach that Course.}
\end{quote}

The \code{if-then} rule of every language says
\code{with(intro: ?vs) \{ \ldots \}}. A reference to a variable of
\code{env.intro} is then a \emph{choice}
\code{\{choice: v, first, later\}} holding both phrases: \code{first} is the
language's wildcard phrase, \code{later} its reference phrase, each in the
slot's case. Which mention is the first is only known when the sentence is
complete, because rules do not build phrases in reading order (an override
is built before the reading it fills, a relative clause before its head may
be chosen, a failed attempt is discarded). Linearisation therefore resolves
choices: it walks the phrase tree in token order and takes \code{first} at
the first choice for a variable and \code{later} at the others. Variables
introduced by their own quantifier (\textit{some Offering}) are not in
\code{intro}; their later mentions are ordinary references.

This is the one place where a rule's decision is carried through the phrase
tree to linearisation; the mechanism is general and could serve other
anaphora (\textit{it}, \textit{they}) later.

% ─────────────────────────────────────────────────────────────────────────────
\section{Linearisation}
\label{sec:linearise}

A phrase is a nested list of tokens \code{\{text, cls\}} with a style class:
\code{kw} (fixed words), \code{type} (object types and variables),
\code{reading}, \code{const} (values) and \code{punct}; libraries may add
classes of their own (German \code{relEnd}). Linearisation (\code{linearise}
in \code{engine.js}):

\begin{enumerate}
  \item flattens the tree in order, resolving choices
        (Section~\ref{sec:choices});
  \item lets the language contract neighbouring tokens
        (\code{lang.contract}): French \textit{que il} $\to$ \textit{qu'il},
        Italian \textit{un studente} $\to$ \textit{uno studente} and
        \textit{in il} $\to$ \textit{nel}, Spanish \textit{y I\ldots} $\to$
        \textit{e I\ldots}, the German closing comma of a relative clause
        dropped before punctuation and at the end, and the English relative
        \textit{that} dropped before a back-reference (\textit{each Course
        that Student enrols in}, not \textit{that that Student});
  \item puts a space before every token except attaching punctuation, a
        token marked \code{glue} and the token after an opening parenthesis;
  \item capitalises the first letter and adds the full stop.
\end{enumerate}

Object type names keep their own capitalisation. The Constraints tab colours
the tokens like the formulas.

% ─────────────────────────────────────────────────────────────────────────────
\section{Population Sentences}
\label{sec:population}

The Population tab's Sentences view writes each fact and instance as a
sentence (\code{population.js}). It does not go through the rules: a
population row is not a formula. It uses the same helpers
(\code{phraseKit(ctx, language)} gives the \code{h} the rules get, so
readings, lexicon, hyphen binding and shaping apply) and the same
linearisation, with fixed words and phrase shapes from each library's
\code{population} section (\textit{the}, \textit{with}, \textit{in which},
\textit{there is}, \textit{is a}, \textit{are referred to by},
\textit{(derived)}).

An instance is a noun phrase found from its type's preferred identifier as
the cell holds it: a popular reference mode in the short form
\textit{Person 'P1'} (with one key line per list, \textit{Persons are
referred to by PersonNr}), a unit-based one \textit{Mass 12 kg}, a general one
\textit{Book with ISBN '\ldots'}, a composite identifier through its
identifying readings (\textit{the Room that is in Building 'MAL' and has
RoomNr '151'}), an objectified fact type through its nominal reading
(\textit{the enrolment of Student 'P1' in Course 'DB1'}) or else
\textit{the Enrolment in which Student 'P1' enrols in Course 'DB1'}. Values
are quoted unless their datatype is a number or a truth value. A callback
lets Dutch ask for the subordinate order and Polish and German give each slot
its case.

% ═════════════════════════════════════════════════════════════════════════════
\part{The Rule Language}
\label{part:dsl}

% ─────────────────────────────────────────────────────────────────────────────
\section{Rule Files}
\label{sec:syntax}

A language's rules are in a rule file (\code{languages/<lang>.vr}), parsed
and compiled by \code{rules.js}. A file holds a language declaration and
rules; \code{\#} starts a comment.

\begin{lstlisting}
language en "English"

sentence rule each-at-most {
  # forall x (exists ys R -> count<=n ys R) says no more than
  # forall x count<=n ys R
  match:    ForAll([?x], Implies(?a, ?q @ Count(<=, ?n, ?ys, ?c)))
  guard:    sameUpToWild(?a, ?c, ?ys)
  template: "each" name(?x) vp(?q, ?x)
}
\end{lstlisting}

A rule is written \code{<kind> rule <name> \{ \ldots\ \}} with the fields \code{match}, \code{guard} and \code{template}, and kind \code{sentence}, \code{clause} or
\code{vp}; the guard is optional. Rules are tried in file order within their
kind. Names may contain hyphens. Errors carry the line and column.

\subsection{Patterns}

Patterns are compiled to data patterns of \code{match.js}; a formula node
matches field by field, in order, so a quantifier's variables are bound
before its body.

\begin{lstlisting}
_                       anything, binds nothing
?x                      binds a value; a second occurrence must be the
                        same value (variables compared by identity)
?xs*                    in a list: the remaining elements
?q @ P                  binds the node P matched to ?q
"text"  12              a string or number
<= >= = .. < > != ~     operators (also as Unicode: <= is the same as the
                        less-or-equal sign, ~ the approximately-equal sign)
[P, ...]                a list of exactly these elements (or a ?tail*)
{ field: P, ... }       a node with these fields
ForAll(vars, body)      Exists(vars, body)    Count(op, n, vars, body)
Not(arg)  And(args)  Or(args)  Implies(l, r)  Iff(l, r)
Type(T, x)  Cmp(op, l, r)  VCmp(op, l, r)  TC(vars, body, args)
Fact(?P, {?x, ?y}, ?rest*)
                        a fact atom over ?P whose arguments include the
                        bound ?x and ?y, in any roles; without ?rest* no
                        other arguments are allowed
Node("dpred")  True     a node of a type; the true node
\end{lstlisting}

Fact arguments are matched by role, not by position. In a vp rule
\code{?subj} is bound to the subject before matching starts. Comparisons of
subformulas (\code{sameUpToWild}, \code{exclusionItems}) use equality up to
renaming of bound variables, because \code{formalise()} builds each copy of a
join path with fresh variables.

\subsection{Guards}

A guard is an expression: bindings \code{?x}, strings, numbers,
\code{true}, \code{false}, \code{null}, lists \code{[a, b]}, field access
\code{?a.t}, \code{env.verb}, calls \code{f(a, b)} of built-ins and library
functions, comparisons \code{==}, \code{!=}, \code{<}, \code{>}, \code{<=},
\code{>=}, membership \code{a in [\ldots]}, and \code{not}, \code{and},
\code{or}. A guard passes when its value is truthy. Guards test structure;
they must not build phrases. Whether a phrase can be built is decided by the
template failing.

\subsection{Templates}

A template is a sequence of items; it fails (gives \code{null}) when an item
does.

\begin{predtable}
  \predrow{"words"}{fixed words (class \code{kw})}
  \predrow{comma colon semicolon}{punctuation, attached to the preceding word}
  \predrow{f(a, b, key: value)}{a call of a built-in or library function; an argument is a sequence of items, passed as is when it is a single item; a failed argument fails the call}
  \predrow{(item item \ldots)}{a group}
  \predrow{if c \{ \ldots \} else \{ \ldots \}}{a choice; \code{else} is optional}
  \predrow{with(key: value) \{ \ldots \}}{the items with a changed environment (Section~\ref{sec:env})}
  \predrow{node}{the matched formula}
\end{predtable}

\subsection{Built-ins}

\begin{predtable}
  \predrow{name(v)}{the variable's name with subscript, styled as a type}
  \predrow{typeName(T)}{a type's name in the language}
  \predrow{ref(v, case: c)}{a reference: \code{refNP}, \code{wildNP} for a wildcard, a choice for a variable of \code{env.intro}}
  \predrow{sentence(F), clause(F, key: v), vp(F, v, key: v)}{a sub-phrase of that kind, with optional environment changes}
  \predrow{reading(A, subject: v)}{an atom through a selected reading}
  \predrow{term(t)}{a data term, through the library's \code{term}}
  \predrow{kw(text), punct(text)}{a fixed word, punctuation}
  \predrow{len(xs), any(xs, p), all(xs, p)}{list tests for guards}
\end{predtable}

Every other name is a library function, called with the phrase helpers
first: \code{lib.f(h, a, b, \{key: value\})}.

% ─────────────────────────────────────────────────────────────────────────────
\section{The Library Interface}
\label{sec:library-interface}

A language is its rule file compiled against its library
(\code{compileLanguage(src, lib)}): \code{\{id, label, rules, \ldots hooks\}}.
The library is a JavaScript module with \code{functions} (called from
guards and templates) and hooks the engine calls:

\begin{predtable}
  \predrow{refNP(h, v, c)}{a reference to an introduced variable in case \code{c}}
  \predrow{wildNP(h, w, c)}{a wildcard (also a variable's first mention)}
  \predrow{reflexive(h, v, c)}{a second occurrence of the subject}
  \predrow{term(h, t)}{a data term or aggregate as a noun phrase}
  \predrow{verbForm(h, text)}{the first verb of a reading in \code{env.verb}'s form, or \code{null}}
  \predrow{pluralise(name)}{a plural by rule, when the lexicon gives none}
  \predrow{shapeReading(h, r, elided, atom)}{a reading's parts for the environment (word order, verb forms), optionally with changed cases}
  \predrow{contract(tokens)}{the token-level rewrites of linearisation}
  \predrow{population}{the fixed words and shapes of population sentences}
\end{predtable}

The phrase helpers \code{h} give the environment (\code{h.env}), the context,
\code{h.with(env)}, sub-phrases (\code{h.sentence}, \code{h.clause},
\code{h.vp}), names (\code{h.name}, \code{h.plural}, \code{h.typeLabel},
\code{h.subscript}), the lexicon entry and gender of a variable's type
(\code{h.entry}, \code{h.gender}), \code{h.ref} and \code{h.reading}.

\paragraph{Shared structural tests.}
\code{structure.js} holds the tests every language uses, so that languages
differ only in wording:

\begin{predtable}
  \predrow{rangeSpec(F, x)}{F is a value range on x: a disjunction of $x \approx c$, $x \ge a$, $x \le b$ and intervals with open or closed bounds}
  \predrow{exclusionItems(F)}{$\neg(A_1 \wedge \ldots \wedge A_n)$, or $\neg(A_i \wedge A_j)$ for every pair of $n \ge 3$ formulas}
  \predrow{exclusiveOrItems(F)}{$(A_1 \vee \ldots \vee A_n)$ together with \textit{at most one of} them}
  \predrow{sameUpToWild(A, B, ys)}{A is B with ys replaced by wildcards}
  \predrow{distributable(ys, args, x)}{$\exists^{=1}\bar y\,\bigwedge R_i(x,y_i)$ with one binary conjunct per $y_i$, which holds exactly when each $\exists^{=1} y_i\,R_i(x,y_i)$ does}
  \predrow{isComplex(F), isSingular(op, n), subjectOf(F), occurs(F, v)}{small tests}
\end{predtable}

\paragraph{Shared phrase builders.}
Each library has its own version of a small set of builders the rules call:
\code{objectQuantified} (a verb phrase from an atom in F with ?y as object
filled by a noun phrase, F's other conjuncts a relative clause on ?y),
\code{anchoredClause} (the same as a clause, the atom anchored at an
introduced variable), \code{relative} and \code{relativeBody} (\textit{that}
+ verb phrase, else an object gap, else \textit{such that} + clause),
\code{quantNP}, \code{clauseList}, \code{names}, \code{commaList}.

% ═════════════════════════════════════════════════════════════════════════════
\part{The Languages}

% ─────────────────────────────────────────────────────────────────────────────
\section{English}
\label{sec:languages}
\label{sec:english}

English follows NORMA's wording. Its 54 rules, in the order they are tried, by phrase kind.

\needspace{12\baselineskip}
\paragraph{Sentence rules}
\begin{center}
\small
\begin{tabularx}{\textwidth}{>{\raggedright\arraybackslash\leavevmode\ttfamily\color{codefg}}p{0.3\textwidth} X}
\toprule
\normalfont\textbf{Rule} & \textbf{Matches $\to$ example} \\
\midrule
impossible, possible & the mode nodes $\to$ \textit{It is impossible that \ldots}, \textit{It is possible that \ldots} \\
obligatory, forbidden & $O(\varphi)$, $F(\chi)$ $\to$ \textit{It is obligatory that \ldots} \\
value-range & $\forall x$ (value range) $\to$ \textit{The possible values of Grade are 0 to 100.} \\
exactly-one-of & exclusive-or $\to$ \textit{For each X, exactly one of the following holds: \ldots} \\
no & $\forall x\,\neg\varphi$ $\to$ \textit{No Person is a Building.} \\
each-at-most & frequency $\to$ \textit{Each Student enrols in at most 6 Courses.} \\
each-that & $\forall x\,(\alpha \to \beta)$ $\to$ \textit{Each Student that enrols in at least 4 Courses is a BusyStudent.} \\
at-least-one-of & inclusive-or $\to$ \textit{For each X, at least one of the following holds: \ldots} \\
each & $\forall x\,\varphi$ $\to$ \textit{Each Person has at most one PersonName.} \\
definition & derivation rules, bare references $\to$ \textit{Student is a TeachingAssistant if and only if Student assists in some Course.} \\
if-then & $\forall\bar x\,(\alpha\to\beta)$, two or more variables, all in $\alpha$; first and later mentions $\to$ \textit{If some Course$_1$ is a prerequisite for some Course$_2$ and \ldots\ then \ldots} \\
for-each & $\forall\bar x\,\varphi$ $\to$ \textit{For each Building and RoomNr, at most one Room is in that Building and has that RoomNr.} \\
clause & anything \\
\bottomrule
\end{tabularx}
\end{center}

\needspace{12\baselineskip}
\paragraph{Clause rules}
\begin{center}
\small
\begin{tabularx}{\textwidth}{>{\raggedright\arraybackslash\leavevmode\ttfamily\color{codefg}}p{0.3\textwidth} X}
\toprule
\normalfont\textbf{Rule} & \textbf{Matches $\to$ example} \\
\midrule
more-than-one-same & negative uniqueness $\to$ \textit{more than one Person has the same Email} \\
at-most-one-of & exclusion $\to$ \textit{at most one of the following holds: A; B} \\
not-subject, none, not & \textit{Course$_1$ is not a prerequisite for Course$_3$}; \textit{no Staff coordinates that Course}; \textit{it is not true that A} \\
and, or, implies, iff, for-each & the connectives \\
some-subject, some-object & $\exists y$ $\to$ \textit{some Enrolment involves \ldots}, \textit{\ldots\ teaches some Offering that \ldots} \\
quantified-subject, -object & singular counts \\
there-are-n-objects, there-are-n-that & cardinality and plural counts $\to$ \textit{there are at least 5 Staff members that coordinate some Course} \\
there-is-some, there-are-n & generic $\exists$ and counts with \textit{such that} \\
fact, subject-vp, value-comparison, data-\-predicate, closure, true, unknown & generic, one per node type, and the catch-all \\
\bottomrule
\end{tabularx}
\end{center}

\needspace{12\baselineskip}
\paragraph{Verb phrase rules}
\begin{center}
\small
\begin{tabularx}{\textwidth}{>{\raggedright\arraybackslash\leavevmode\ttfamily\color{codefg}}p{0.3\textwidth} X}
\toprule
\normalfont\textbf{Rule} & \textbf{Matches $\to$ example} \\
\midrule
fact, is-a, value-comparison, same-object, data-predicate & \textit{has some PersonName}, \textit{is a Student}, \textit{is at least 1970}, \ldots \\
directly-or-indirectly & $[\mathsf{TC}_{v,w}R(v,w)](x,t)$ $\to$ \textit{is a prerequisite for itself, directly or indirectly} \\
and, or, exactly-one-each & joined verb phrases; \textit{is in exactly one Building and has exactly one RoomNr} \\
quantified-object, no-object, some-object & \textit{has at most one Email}, \textit{enrols in no Course} \\
not, none & negated verb phrases (negative mode), \textit{earned no Grade} \\
\bottomrule
\bottomrule
\end{tabularx}
\end{center}

The library (\code{english.js}) holds a small morphology: \textit{one} for 1,
plurals by rule keeping a subscript (\textit{Category} $\to$
\textit{Categories}), \textit{a}/\textit{an}, \textit{who} for personal
types, \textit{itself}/\textit{themselves}, verb inflection for plural and
negated subjects (\textit{is}/\textit{are}/\textit{is not},
\textit{has}/\textit{have}/\textit{does not have}), the wording of data
operators and aggregates (\textit{the sum of the Credits that each Course
that Student enrols in is worth}), and the contraction that drops a relative
\textit{that} before a back-reference.

% ─────────────────────────────────────────────────────────────────────────────
\section{The Other Languages}

Each language has a rule file of 53--57 rules that mirrors the English one,
so that rule names, coverage and golden files line up, and a library that
supplies what the language needs beyond English.

\subsection{Dutch}

Word order is an environment key, \code{env.form}: main clause (verb
second), \code{sub} (verb last) and \code{inv} (verb first). Readings are
written in main-clause order; the subordinate form comes from
\code{subParts}, else the finite verb, the first word after the subject,
moves to the end (\textit{Vak is Studiepunten waard} $\to$ \textit{\ldots\
Vak Studiepunten waard is}), and inversion puts it before the subject. The
\code{if-then} rule therefore reads \code{"als" clause(?a, form: "sub") comma
"dan" clause(?c, form: "inv")}, with an \code{if-then-holds} fallback
\textit{\ldots, dan geldt dat \ldots} for consequents without an inverted
form. Gender (\textit{de}/\textit{het}) chooses \textit{elke}/\textit{elk}
and \textit{die}/\textit{dat}:

\begin{quote}
\textit{Als een Medewerker een Vakaanbieding die een Vak betreft doceert,
dan is die Medewerker bevoegd om dat Vak te doceren.}
\end{quote}

\subsection{French, Italian and Spanish}

The Romance languages share one library (\code{romance.js}) parametrised by
a profile per language (\code{french.js}, \code{italian.js} and \code{spanish.js}). Word order
is the English one. Gender (\textit{m}/\textit{f}) chooses articles and
demonstratives (\textit{un}/\textit{une}, \textit{ce}/\textit{cette},
\textit{quel}/\textit{quella}); negation is a particle before the verb and in
French \textit{pas} after it (\code{env.verb} \code{negated}, or \code{ne}
alone before \textit{aucun}); a preposition before a relative gap moves to
the pronoun (\textit{dans laquelle}, \textit{in cui}, \textit{en la que});
and the profile's contractions run at linearisation (elision, \textit{uno}
before \textit{s}+consonant, Italian articulated prepositions, Spanish
\textit{e} before \textit{i}).

\begin{quote}
\textit{Si un Membre du personnel enseigne une Offre de cours qui concerne
un Cours, alors ce Membre du personnel est qualifié pour enseigner ce
Cours.}
\end{quote}

\subsection{Polish}

Nouns decline. The lexicon gives a type its nominative singular and plural,
gender (\textit{mp}, \textit{ma}, \textit{mi}, \textit{f}, \textit{n}) and its
other case forms; a reading gives each slot a case (nominative for the
first, accusative for the others by default); a negated verb puts an
accusative object in the genitive. Because noun phrases depend on their
slot, the rules pass the \emph{kind} of a noun phrase (\code{"some"},
\code{"none"}, \code{"each"}, \code{"bare"} or a counting quantifier) to
\code{np}, \code{anchoredClause} and \code{objectQuantified}, and the library
builds it in the slot's case. Numbers govern the noun (\textit{jeden
Przedmiot}, \textit{2 Przedmioty}, \textit{5 Przedmiotów}).

\begin{quote}
\textit{Jeśli jakiś Pracownik prowadzi jakąś Ofertę, która dotyczy jakiegoś
Przedmiotu, to ten Pracownik ma uprawnienia do nauczania tego Przedmiotu.}
\end{quote}

\subsection{German}

German combines Dutch word order (verb second, verb last in subordinate
clauses, a separable particle joining the verb at the end:
\textit{teilnimmt}) with Polish-style cases (nominative, genitive, dative,
accusative per slot; noun phrase kinds \textit{jeder}, \textit{ein},
\textit{kein}, \textit{dieser}, \textit{derselbe}, counts). The lexicon
gives gender (else guessed from the ending), plural (else by rule) and case
forms (else by rule: genitive \textit{-s}/\textit{-es}, weak masculine nouns
\textit{-en}, dative plural \textit{-n}). Subordinate clauses take commas on
both sides; the closing one is a \code{relEnd} token that contraction drops
before punctuation and at the end. A bare name reads as a label and stays in
the nominative.

\begin{quote}
\textit{Wenn ein Mitarbeiter ein Kursangebot, das einen Kurs betrifft,
unterrichtet, dann darf dieser Mitarbeiter diesen Kurs unterrichten.}
\end{quote}

% ═════════════════════════════════════════════════════════════════════════════
\part{Tools and Implementation}

% ─────────────────────────────────────────────────────────────────────────────
\section{Tools}
\label{sec:tools}

\subsection{Golden tests}

\code{npm run test:verbalisation} (\code{scripts/verbalisation-golden.mjs})
verbalises every entry of the reference model (the university sample, whole
schema with the implicit entries: 153 entries) and of the additions in
\code{tests/verbalisation/extras.mjs} (cardinality ranges, derived and
semi-derived types and fact types, data terms, aggregates, deontic
constraints, open bounds, an outer join: 29 entries), in all seven languages
and in the positive and negative modes, and compares them with
\code{tests/verbalisation/<sample|extras>.<lang>[.neg].txt}. Each block holds
the entry's kind, its formula and its sentence. \code{--update} accepts new
output, so the diff of a rule change shows exactly which sentences it
affects; \code{--lang=xx} checks one language.

\subsection{Coverage and overlaps}

\code{npm run verbalisation:coverage} (and the workbench) report per
language which rules the reference sentences use, which never fire, and
which sentences fell back to a generic rule (\code{analysis.js}, from the
\code{stats} that \code{verbalise} collects). Overlaps are found statically:
a rule is overlapped by an earlier one of the same kind whose pattern
subsumes its own and whose guard is absent or the same, so it is only tried
when the earlier template fails.

\subsection{Workbench}

View \textrightarrow\ Verbalisation workbench (Cmd+Shift+B) opens a window with the rule file of a
language in an editor, and next to it the entries of the open model with
formula, sentence and trace, in the Sentences, Coverage and Overlaps tabs.
Applied rules are kept per language in local storage and used by the
Constraints tab of the main window; the windows exchange the model and the
rules through the main process (\code{workbench.js}).

% ─────────────────────────────────────────────────────────────────────────────
\section{Implementation}

\subsection{Files}

\begin{predtable}
  \predrow{src/\allowbreak logic/\allowbreak constraintFormulas.js}{stage 0: \code{formalise}, \code{makeContext}, default entries}
  \predrow{src/\allowbreak verbalisation/\allowbreak negative.js}{the negative rewrite}
  \predrow{src/\allowbreak verbalisation/\allowbreak normalise.js}{the wildcard rewrite}
  \predrow{src/\allowbreak verbalisation/\allowbreak engine.js}{\code{verbalise}, rule application, helpers, reading selection, choices, linearisation}
  \predrow{src/\allowbreak verbalisation/\allowbreak match.js}{pattern matcher, equality up to bound variables}
  \predrow{src/\allowbreak verbalisation/\allowbreak rules.js}{the rule language: parser and compiler}
  \predrow{src/\allowbreak verbalisation/\allowbreak structure.js}{shared structural tests}
  \predrow{src/\allowbreak verbalisation/\allowbreak lexicon.js, nameStyle.js, hyphenBinding.js}{lexicon lookup, name display, hyphen binding}
  \predrow{src/\allowbreak verbalisation/\allowbreak population.js}{population sentences}
  \predrow{src/\allowbreak verbalisation/\allowbreak analysis.js, workbench.js}{coverage, overlaps, the workbench's shared state}
  \predrow{src/\allowbreak verbalisation/\allowbreak languages/}{\code{<lang>.vr} and \code{<lang>.js} per language, \code{romance.js}, \code{index.js}}
\end{predtable}

\subsection{Interpretation, not compilation}

Rules are interpreted. Each parsed pattern, guard and template becomes a
closure once, when the rule file is loaded; no JavaScript source is
generated and no user-written code is evaluated in the renderer. Verbalising
the whole reference model takes milliseconds, so the workbench can
re-verbalise everything on every change.

% ─────────────────────────────────────────────────────────────────────────────
\section{Open Questions}
\label{sec:open}

\begin{itemize}
  \item \textbf{Remaining \textit{that that}.} After the conjunction
        \textit{that}, a back-reference still gives \textit{such that that
        Student} and \textit{it is not true that that Student} in negative
        sentences, as in NORMA. Dropping the conjunction would lose the
        reference; rephrasing (\textit{such that the Student}) would give up
        NORMA's consistent \textit{that}.
  \item \textbf{Pronouns.} The choice mechanism (Section~\ref{sec:choices})
        could also write \textit{it} or \textit{they} for a third mention, or
        for a variable mentioned once before in the same clause.
  \item \textbf{First mentions under negation.} A variable of \code{intro}
        whose first mention lies under \textit{it is not true that} would be
        introduced there; no current formula does this, but the rule could
        guard against it.
  \item \textbf{Contraction in the DSL.} Token rewrites (\code{contract})
        are written in the libraries. A rule kind for them would let rule
        writers add elisions and the like without JavaScript.
  \item \textbf{Language files with \code{extends}.} v8's dialect
        inheritance is not implemented; the Romance languages share a
        library instead.
  \item \textbf{Subjunctive.} The negative mode uses the indicative in
        French, Italian and Spanish (\textit{il est impossible que \ldots\
        suit}).
\end{itemize}

% ═════════════════════════════════════════════════════════════════════════════
\appendix

\section{Worked Examples}

Each example gives the formula as the formula view shows it, the sentence,
and the rule trace.

\paragraph{Uniqueness without a reading from the constrained role.}
\[
\forall \mathit{Course}\;\exists^{\le 1}\mathit{Staff}\;(\mathit{Staff}\ \textsf{coordinates}\ \mathit{Course})
\]
\textit{For each Course, at most one Staff member coordinates that Course.}\\
Trace: \code{sentence:for-each} $\to$ \code{clause:quantified-subject}
$\to$ \code{vp:fact}. The \code{each} rule failed first, because
\textit{coordinates} has no reading starting with Course.

\paragraph{Frequency.}
\[
\forall \mathit{Student}\,((\exists \mathit{Course}\,(\mathit{Student}\ \textsf{enrols in}\ \mathit{Course})) \to \exists^{\le 6}\mathit{Course}\,(\mathit{Student}\ \textsf{enrols in}\ \mathit{Course}))
\]
\textit{Each Student enrols in at most 6 Courses.}\\
Trace: \code{sentence:each-at-most} $\to$ \code{vp:quantified-object}.

\paragraph{Subset over a join path, positive.}
\[
\begin{aligned}\forall \mathit{Staff}\,\forall \mathit{Course}\,(&(\exists \mathit{Offering}\,((\mathit{Staff}\ \textsf{teaches}\ \mathit{Offering}) \wedge (\mathit{Offering}\ \textsf{involves}\ \mathit{Course})))\\ &\to (\mathit{Staff}\ \textsf{is qualified to teach}\ \mathit{Course}))\end{aligned}
\]
\textit{If some Staff member teaches some Offering that involves some Course
then that Staff member is qualified to teach that Course.}\\
Trace: \code{sentence:if-then} $\to$ \code{clause:some-object} $\to$
\code{vp:fact} $\to$ \code{clause:fact}. \textit{Staff member} and
\textit{Course} are in \code{intro}; \textit{Offering}, used twice, keeps its
own quantifier.

\paragraph{The same, negative.}
\[
\neg(\exists \mathit{Staff}\,\exists \mathit{Course}\,((\exists \mathit{Offering}\,(\ldots)) \wedge \neg(\mathit{Staff}\ \textsf{is qualified to teach}\ \mathit{Course})))
\]
\textit{It is impossible that there is some Staff member and some Course
such that that Staff member teaches some Offering that involves that Course
and that Staff member is not qualified to teach that Course.}

\paragraph{Acyclic ring.}
\[
\forall \mathit{Course}_1\;\neg[\mathsf{TC}_{\mathit{Course}_2,\mathit{Course}_3}\,(\mathit{Course}_2\ \textsf{is a prerequisite for}\ \mathit{Course}_3)](\mathit{Course}_1,\mathit{Course}_1)
\]
\textit{No Course$_1$ is a prerequisite for itself, directly or indirectly.}\\
Trace: \code{sentence:no} $\to$ \code{vp:directly-or-indirectly}.

\paragraph{External preferred identifier.}
\[
\forall \mathit{Room}\;\exists^{=1}\mathit{Building},\mathit{RoomNr}\,((\mathit{Room}\ \textsf{is in}\ \mathit{Building}) \wedge (\mathit{Room}\ \textsf{has}\ \mathit{RoomNr}))
\]
\textit{Each Room is in exactly one Building and has exactly one RoomNr.}\\
Trace: \code{sentence:each} $\to$ \code{vp:exactly-one-each} $\to$
\code{vp:quantified-object} (twice).

\paragraph{Default verbalisation.}
\[
\Diamond\,\exists \mathit{PersonName}\;\exists^{>1}\mathit{Person}\,(\mathit{Person}\ \textsf{has}\ \mathit{PersonName})
\]
\textit{It is possible that more than one Person has the same PersonName.}

\paragraph{Deontic, negative.}
\[
F(\exists \mathit{Score}\,((\mathit{Score} \le 0 \vee \mathit{Score} > 100) \wedge \mathit{Score} \le 200))
\]
\textit{It is forbidden that some Score is at most 0 or is more than 100 and
is at most 200.}

% ─────────────────────────────────────────────────────────────────────────────
\section{Changes from v9}

\begin{itemize}
  \item The document covers the whole pipeline, with the rule language as
        Part~II and the languages as Part~III.
  \item Rules are written in \code{.vr} files with the concrete syntax of
        Section~\ref{sec:syntax}, parsed and compiled by \code{rules.js};
        each language is a rule file plus a library with a fixed hook
        interface (Section~\ref{sec:library-interface}).
  \item Seven languages: English, Dutch, French, Italian, Spanish, Polish and
        German, with a per-language lexicon in the model (names, plurals,
        genders, case forms; readings with subordinate and plural forms and
        cases) and a view language.
  \item Presentation modes: negative (a counterexample rewrite), default
        ($\Diamond$ entries from \code{formalise()}) and deontic
        ($O$, $F$).
  \item New environment keys \code{verb}, \code{same}, \code{none},
        \code{form} and \code{intro}; \textit{If \ldots\ then \ldots}
        sentences introduce their variables with \textit{some} and refer
        back with \textit{that}, resolved at linearisation as choices.
  \item Linearisation contracts tokens per language.
  \item Hyphen binding, name display styles and population sentences.
  \item Coverage and overlap reports and the workbench window exist.
\end{itemize}

\end{document}
